%%%%%%%%%%%%%%%%%%%%%%%%%%%%%%%%%%%%%%%%%%%%%%%%%%%%%%%%%%%%%%%%%%%%%%
%                                                                    %
% How to define the Riesz transform for Fourier-Neumann expansions?  %
%                                                                    %
% CONTRIBUCIONES CIENT\'IFICAS EN HONOR DE MIRIAN ANDR\'ES G\'OMEZ,  %
% Laureano Lamb\'an, Ana Romero y Julio Rubio (Editores),            %
% Servicio de Publicaciones, Universidad de La Rioja.                %
% Logro\~no, 2010                                                    %
% ISBN 978-84-96487-50-5                                             %
%                                                                    %
% DATE: first version, 2-February-2004                               %
% DATE: last modification, 29-March-2010                             %
%                                                                    %
% Written with AMS-LaTeX 2.0                                         %
%                                                                    %
%%%%%%%%%%%%%%%%%%%%%%%%%%%%%%%%%%%%%%%%%%%%%%%%%%%%%%%%%%%%%%%%%%%%%%

\documentclass{amsart}
%\documentclass[11pt,a4paper]{amsart}

\usepackage[twoside, 
paperwidth=17cm, paperheight=24cm,
textheight=19cm, textwidth=12.5cm,
top=3.0cm, centering]{geometry}

\usepackage[spanish, english]{babel}

%% S\'olo se usa si el art\'{\i}culo est\'a en espa\~nol:
%\input{libroMirianAndres.sty}
%% Si est\'a en ingl\'es basta esto:
\def\publname{CONTRIBUCIONES CIENT\'IFICAS\\ 
EN HONOR DE MIRIAN ANDR\'ES G\'OMEZ\\
(Laureano Lamb\'an, Ana Romero y Julio Rubio, editores),\\
Servicio de Publicaciones, Universidad de La Rioja,\\
Logro\~no, Spain, 2010.}
%%% CON ESTO LO LOGRO QUITAR TODO SALVO EL \publname
\copyrightinfo{}{} %%% {copyright year}{copyright holder}
\issueinfo{}{}{}{} %%% {vol}{number}{month}{year}
\PII{\null}        %%% ISSN and reference% (Sin el \null, da siempre error de underfull)
\setcounter{page}{203}

%\usepackage[latin1]{inputenc}

%\usepackage{amsmath,amsthm} %%% amsthm necesario para usar el entorno proof

\usepackage{amssymb}

%%% FOR DRAFTS:
% \overfullrule=5pt
% \usepackage[notcite,notref]{showkeys} 
%\date{August 20, 2009}

%%%%%%%% A LOT OF BIG SEVERAL-LINE EQUATIONS (aligns, etc.).
%%%%%%%% THUS, WE ALLOW TO BREAK THEM IN PAGES.
%\allowdisplaybreaks

%\theoremstyle{definition}
%\newtheorem*{rem*}{Remark}

\newcommand{\spa}{\operatornamewithlimits{span}} %%%% \span already exists in TeX
%\newcommand{\supp}{\operatornamewithlimits{supp}}
%\newcommand{\iden}{\operatorname{Id}}
%\renewcommand{\Re}{\operatorname{Re}}
%\renewcommand{\Im}{\operatorname{Im}}

%%%%%%%%%%%%%%%%%%%%%%%%%%%%%%%%%%%%%%%%%%%%
\begin{document}
%%%%%%%%%%%%%%%%%%%%%%%%%%%%%%%%%%%%%%%%%%%%

%%%%%%%%%%%%%%%%%%%%%%%%%%%%%%%%%%%%%%%%%%%%%%
\title[How to define the Riesz transform?]{How to define the Riesz transform for Fourier-Neumann expansions?}
%\author[\'O. Ciaurri, L. Roncal, J. L. Torrea, and J. L. Varona]%
%  {\'Oscar Ciaurri, Luz Roncal, Jos\'e L. Torrea, and Juan L. Varona}

\author[\'O. Ciaurri]{\'Oscar Ciaurri}
\address{Departamento de Matem\'aticas y Computaci\'on \\
         Universidad de La Rioja \\
         26004 Logro\~no, Spain}
\email{oscar.ciaurri@unirioja.es}

\author[L. Roncal]{Luz Roncal}
\address{Departamento de Matem\'aticas \\
         IES Hermanos d'Elhuyar \\
         26003 Logro\~no, Spain}
\email{luz.roncal@unirioja.es}

\author[J. L. Torrea]{Jos\'e L. Torrea}
\address{Departamento de Matem\'aticas \\
         Facultad de Ciencias \\
         Universidad Aut\'o\-no\-ma de Madrid \\
         28049 Madrid, Spain}
\email{joseluis.torrea@uam.es}
\urladdr{http://www.uam.es/joseluis.torrea}

\author[J. L. Varona]{Juan L. Varona}
\address{Departamento de Matem\'aticas y Computaci\'on \\
         Universidad de La Rioja \\
         26004 Logro\~no, Spain}
\email{jvarona@unirioja.es}
\urladdr{http://www.unirioja.es/cu/jvarona}

\thanks{Research supported by grants 
MTM2009-12740-C03-03 and % <- La Rioja \'ultimo 
MTM2008-06621-C02-01 of the~DGI} % <- Torrea \'ultimo 

\keywords{Fourier-Neumann expansions, Heat semigroup,
Poisson semigroup, Conjugate Poisson semigroup,
Fractional integrals, Riesz potentials, Riesz transform.}

\subjclass[2000]{Primary 42C10}
%\subjclass[2000]{Primary 42C10; Secondary 35K05, 35J05}
%%%%%%%%%%%%%%%%%%%%%%%%%%%%%%%%%%%%%%%%%%%%%%

\dedicatory{Dedicated to the memory of our colleague Mirian}

%\renewcommand{\abstractname}{Resumen}

%%% ABSTRACT EN DOS IDIOMAS CON amsart

%% Primero ingl\'es y luego espa\~nol:
%\selectlanguage{english}
%\begin{abstract}
%This is the abstract in English. ......
%...... ...... ...... ...... ...... ......
%...... ...... ...... ...... ...... ......
%...... ...... ...... ...... ...... ......
%...... ...... ...... ...... ...... ......
%\medskip
%\selectlanguage{spanish}
%\item[\hskip\labelsep\scshape\abstractname.]
%Este es el abstract en espa\~nol. ......
%...... ...... ...... ...... ...... ......
%...... ...... ...... ...... ...... ......
%...... ...... ...... ...... ...... ......
%...... ...... ...... ...... ...... ......
%\end{abstract}
%\selectlanguage{english} % or \selectlanguage{spanish}

% Primero espa\~nol y luego ingl\'es:
\selectlanguage{spanish}
\begin{abstract}
Consideramos un operador diferencial autoadjunto $L$ con autovalores positivos.
En este art\'{\i}culo damos una visi\'on general del m\'etodo 
que se sigue para definir de manera espectral el semigrupo del calor $e^{-tL}$, 
el de Poisson $e^{-t\sqrt{L}}$ y los potenciales de Riesz $L^{-a}$, 
hasta llegar a la descomposici\'on $L = \delta^*\delta$ (con $\delta$
y $\delta^*$ dos operadores diferenciales de primer orden adjuntos),
y as\'{\i} definir la transformada de Riesz $R = \delta L^{-1/2}$.

En particular, para el operador diferencial $L_\alpha$
que se corresponde con los desarrollos de Fourier-Neumann de orden $\alpha>-1$,
con la ayuda de un sistema de \'algebra computacional hemos encontrado
una pl\'eyade de descomposiciones $L_\alpha = \delta^* \delta$ (o similares).
Pero ninguna de ellas parece ser \'util para definir la transformada de Riesz,
puesto que los operadores $\delta$ y $\delta^*$ son demasiado ``raros'': 
% las funciones presentes en las expresiones de $\delta$ y $\delta^*$ 
% tienen una cantidad infinita de polos, luego al aplicar $\delta$ no  
% obtenemos funciones integrables.
las expresiones de $\delta$ y $\delta^*$ contienen funciones con
una cantidad infinita de polos, luego al aplicar $\delta$ no  
obtenemos funciones integrables.
\medskip
\selectlanguage{english}
\item[\hskip\labelsep\scshape\abstractname.]
We consider a self-adjoint differential operator $L$ with positive eigenvalues. 
In this paper, we provide a general view of the method followed to define, 
in a spectral way, the heat semigroup $e^{-tL}$, 
the Poisson semigroup $e^{-t\sqrt{L}}$ and the Riesz potentials $L^{-a}$, 
up to arrive at the decomposition $L = \delta^*\delta$ (with $\delta$
and $\delta^*$ two adjoint first order differential operators),
and then define the Riesz transform $R = \delta L^{-1/2}$.

In particular, for the differential operator  $L_\alpha$
corresponding with the Fourier-Neumann expansions of order $\alpha>-1$,
we have found, with the help of a computer algebra system, 
a pleyade of strange decompositions $L_\alpha = \delta^* \delta$ (or similar).
Nevertheless, none of them seems to be useful to define the Riesz transform,
because operators $\delta$ and $\delta^*$ are too ``weird'': 
% the functions in the expressions of $\delta$ and $\delta^*$ have 
% an infinite quantity of poles, so we do not obtain integrable functions
% by applying~$\delta$.
the expressions of $\delta$ and $\delta^*$ have functions with
an infinite quantity of poles, so we do not obtain integrable functions
by applying~$\delta$.
\end{abstract}
\selectlanguage{english} % or \selectlanguage{spanish}


%%%%%%%%%%%%%%%%%%%%%%%%%%%%%%%%%%%%%%%%%%%%%%
%\selectlanguage{spanish}
\maketitle
%%%%%%%%%%%%%%%%%%%%%%%%%%%%%%%%%%%%%%%%%%%%%%


%%%%%%%%%%%%%%%%%%%%%%%%%%%%%%%%%%%%%%%%%%%%%%%%%%%%%%
\section{A sketch of the general theory}
%%%%%%%%%%%%%%%%%%%%%%%%%%%%%%%%%%%%%%%%%%%%%%%%%%%%%%

Let $X$ be an open (possibly unbounded) interval on the real line,
a weight $w \in C^1(X)$ with $w(x) > 0$ for all $x \in X$,
and let $\langle\cdot,\cdot\rangle$ denote the natural inner
product in $L^2(X,w(x)\,dx)$.

Moreover, let $L$ be a self-adjoint second order differential operator
defined initally on a suitable (and dense) subset of functions
in $L^2(X,w(x)\,dx)$.

Let us assume that there is an orthonormal basis $\{\varphi_n\}_{n=0}^{\infty}$
in $L^2(X,w(x)\,dx)$ consisting of eigenfunctions of $L$ that correspond
with the eigenvalues $\{\lambda_n\}_{n=0}^{\infty}$, that is,
$$
  L \varphi_n = \lambda_n \varphi_n,
$$
and satisfying $0 < \lambda_0 < \lambda_1 < \lambda_2 < \cdots$ and
$\lim_n \lambda_n = +\infty$.

For a function $f \in L^2(X,w(x)\,dx)$ whose Fourier expansion is
$$
  f = \sum_{n=0}^{\infty} c_n \varphi_n,
  \quad c_n = \langle f,\varphi_n \rangle
  = \int_{X} f(x) \overline{\varphi_n(x)} w(x)\,dx,
$$
we define in a natural way (at least formally)
$$
  L f = L\left(\sum_{n=0}^{\infty} c_n \varphi_n\right)
  = \sum_{n=0}^{\infty} c_n \lambda_n \varphi_n.
$$
This spectral method can be used to extend to $L^2(X,w(x)\,dx)$
many operators related to~$L$. For instance, for a function $\Phi$
defined over the eingenvalues we would define the operator $\Phi(L)$
as
$$
  \Phi(L) f = \Phi(L) \left(\sum_{n=0}^{\infty} c_n \varphi_n\right)
  = \sum_{n=0}^{\infty} c_n \Phi(\lambda_n) \varphi_n.
$$

In particular, for $t>0$, we define the operator $e^{-tL}$ by means of
$$
  e^{-tL} \varphi_n = e^{-t\lambda_n} \varphi_n
$$
and, for $f = \sum_{n=0}^{\infty} c_n \varphi_n \in L^2(X,w(x)\,dx)$,
$$
  (e^{-tL} f)(x) = e^{-tL}\left(\sum_{n=0}^{\infty} c_n \varphi_n\right)(x)
  = \sum_{n=0}^{\infty} c_ne^{-t\lambda_n} \varphi_n(x).
$$
(Let us notice that, here and in what follows, the derivatives of the
differential operator $L$ are taken with respect to the variable~$x$.)
If we denote $u(x,t) = (e^{-tL} f)(x)$, the partial derivative with respect
to $t$ is, formally,
$$
  \frac{\partial}{\partial t} u(x,t) = -L e^{-tL} f(x) = -L u(x,t)
$$
so $u(x,t)$ is the solution of the initial value problem
\begin{equation}
\label{eq:heat}
  \left\{
  \begin{aligned}
    &\frac{\partial}{\partial t} u(x,t) = -L u(x,t), \quad x \in X,\ t>0, \\
    &\lim_{t\to0} u(x,t) = u(x,0) = f(x).
  \end{aligned}
  \right.
\end{equation}
In the classical case of $L$ being the ``ordinary'' Laplacian
(with a change of sign to get a differential operator with positive eigenvalues)
$L = - \frac{\partial^2}{\partial x^2}$,
\eqref{eq:heat}~is the heat (or heat-diffusion) equation
which describes the distribution of heat (i.e., variation in temperature)
in a given interval $X$ over time\footnote{The heat equation is derived from Fourier's law and conservation of energy. The first solution technique for the heat equation was proposed by Joseph Fourier in his treatise \textit{Th\'eorie Analytique de la Chaleur}, published in 1822. Later, the solution technique was extended to many other types of equations and this led naturally to the basic ideas of the spectral theory.}
$$
  \left\{
  \begin{aligned}
    &\frac{\partial}{\partial t} u(x,t) = \frac{\partial^2}{\partial x^2} u(x,t), \quad x \in X,\ t>0 ,\\
    &u(x,0) = f(x), \quad x \in X.
  \end{aligned}
  \right.
$$
By the reason above, and since $e^{-t_1L}e^{-t_2L}f = e^{-(t_1+t_2)L}f$ for $t_1,t_2\ge0$,
the operator $e^{-tL}$ is known as the heat semigroup.

Another useful operator is $e^{-t\sqrt{L}}$, that is defined by means of
$$
  e^{-t\sqrt{L}} \varphi_n = e^{-t\sqrt{\lambda_n}} \varphi_n
$$
and, for $f = \sum_{n=0}^{\infty} c_n \varphi_n \in L^2(X,w(x)\,dx)$,
$$
  (e^{-t\sqrt{L}} f)(x) = e^{-t\sqrt{L}}\left(\sum_{n=0}^{\infty} c_n \varphi_n\right)(x)
  = \sum_{n=0}^{\infty} c_ne^{-t\sqrt{\lambda_n}} \varphi_n(x).
$$
This time, if we take $u(x,t) = (e^{-t\sqrt{L}} f)(x)$, we have
$$
  \frac{\partial^2}{\partial t^2} u(x,t) = L e^{-t\sqrt{L}} f(x) = L u(x,t),
$$
so $u(x,t)$ is the solution of
\begin{equation}
\label{eq:Poisson}
  \left\{
  \begin{aligned}
    &\frac{\partial^2}{\partial t^2} u(x,t) = L u(x,t), \quad x \in X,\ t>0, \\
    &\lim_{t\to0} u(x,t) = u(x,0) = f(x), \quad x \in X.
  \end{aligned}
  \right.
\end{equation}
In the classical case of $L = - \frac{\partial^2}{\partial x^2}$ the first
equation in~\eqref{eq:Poisson} becomes
$$
  \Bigg(\frac{\partial^2}{\partial x^2} + \frac{\partial^2}{\partial t^2} \Bigg) u(x,t) = 0,
$$
so it means that $u(x,t)$ is harmonic (in the classical sense).
In general, $e^{-t\sqrt{L}}$ is known as the Poisson semigroup, and a function $u(x,t)$
that satisfies
\begin{equation}
\label{eq:harmonic}
  \Bigg(\frac{\partial^2}{\partial t^2} - L \Bigg) u(x,t) = 0
\end{equation}
is said to be harmonic.

Under some conditions that we will not detail, the formal definitions that we have
shown above can be rigorized on $L^2(X,w(x)\,dx)$ in a standard way; see~\cite{Stein} for details.
Much more difficult is to extend it to functions on $L^p(X,w(x)\,dx)$ for $p \ne 2$.
Several cases of this topic have been studied for different operators and orthogonal systems,
starting from the seminal papers of Muckenhoupt and Stein in 1965~\cite{MuSt} and Muckenhoupt
in 1969~\cite{Mu}. In the recent years a considerable activity can be observed; 
see~\cite{BCiMPTV, MST, Sj, StT, Th1, Th2, Th3, T}.
% the contributions of Thangavelu (such as~\cite{Th1, Th2, Th3}),
% and many papers published by Sj\"ogren, Stempak, Torrea and other collaborators
% (for instance, \cite{BCiMPTV, MST, Sj, StT, T}).
An important part of this work is to identify the operators by means of suitable
integral expressions with kernels, and to find precise bounds for these kernels.

There is a way of relating the Poisson and heat semigroup. For that purpose, let us note the
well known identity
\[
   e^{-t\sqrt{\gamma}} =
   \frac{1}{\sqrt{4\pi}} \int_0^\infty t e^{-t^2/(4s)}e^{-s\gamma} s^{-3/2}\,ds.
\]
%% HECHO CON MATHEMATICA:
% (1/Sqrt[4 Pi])*
% Integrate[t*E^(-t^2/(4 s)) E^(-g*s) s^(-3/2), {s, 0, Infinity},
%  Assumptions -> {t > 0, g > 0}]
% da
% E^(-Sqrt[g] t)
Then, as shown in~\cite{Stein}, the Poisson and heat semigroup can be related
by means of the subordination formula
%\[
%  P_tf(x)
%  = \frac{1}{\sqrt{4\pi}} \int_0^\infty t e^{-t^2/(4s)} W_sf(x) s^{-3/2}\,ds.
%\]
\[
  (e^{-t\sqrt{L}}f)(x)
  = \frac{1}{\sqrt{4\pi}} \int_0^\infty t e^{-t^2/(4s)} (e^{-sL}f)(x) s^{-3/2}\,ds.
\]

There are other operators that can be defined by this technique.
In particular, the formula
%$s^{-\lambda} = \frac{1}{\Gamma(\lambda)} \int_0^\infty t^{\lambda-1} e^{-ts}\,dt$
$$
  s^{-a} = \frac{1}{\Gamma(a)} \int_0^{\infty} t^{a-1} e^{-ts} \,dt
$$
suggests the definition of the Riesz potentials (also known as fractional integrals) either  as
\begin{equation*}
%\label{eq:frac}
%  L^{-\lambda}f(x) =
%  \frac{1}{\Gamma(\lambda)}\int_0^\infty t^{\lambda-1} W_tf(x)\,dt
  L^{-a}f(x) =
  \frac{1}{\Gamma(a)}\int_0^\infty t^{a-1} (e^{-tL}f)(x)\,dt
\end{equation*}
or
\begin{equation*}
%\label{eq:frac/2}
%   L^{-\lambda/2}f(x)
%   = \frac{1}{\Gamma(\lambda)} \int_0^\infty t^{\lambda-1} P_tf(x) \,dt.
   L^{-a/2}f(x)
   = \frac{1}{\Gamma(a)} \int_0^\infty t^{a-1} (e^{-t\sqrt{L}}f)(x) \,dt,
\end{equation*}
that, in a espectral way (for $f = \sum_{n=0}^{\infty} c_n \varphi_n \in L^2(X,w(x)\,dx)$), is
\begin{equation*}
  L^{-a}f(x) = \sum_{n=0}^{\infty} c_n \lambda_n^{-a}  \varphi_n.
\end{equation*}

%%%%%%%%%%%%%%%%%%%%%%%%%%%%%%%%%%%%%%%%%%%%%%%%%%%%%%
\subsection{The decomposition of the differential operator}
%%%%%%%%%%%%%%%%%%%%%%%%%%%%%%%%%%%%%%%%%%%%%%%%%%%%%%

Now, let us suppose that we have a decomposition
\begin{equation}
\label{eq:decomposition}
  L = \delta^* \delta
\end{equation}
where $\delta$ and $\delta^*$ are two adjoint first order differential operators, i.e., that verify
\begin{equation}
\label{eq:adjoint}
  \int_X f(x) \delta g(x) w(x)\,dx
  = \int_X g(x)\delta^*f(x) w(x)\,dx
\end{equation}
for functions $f$ and $g$ good enough.
(In the classical case $L= -\frac{\partial^2}{\partial x^2}$, we have $\delta^* = -\frac{\partial}{\partial x}$,
$\delta = \frac{\partial}{\partial x}$, and \eqref{eq:adjoint} is just the formula of integrating by parts for functions $f,g$ that vanishes in both extremes of the interval~$X$.)
This decomposition $L = \delta^* \delta$ is the key to generalize many classical operators of the harmonic analysis under this general scheme.

In particular, the operator $\delta$ in this decomposition is used to define the so-called Riesz transform
$$
  R = \delta L^{-1/2}
$$
that, again in a spectral way (recall that $\{\varphi_n\}_{n=0}^{\infty}$ is an orthonormal basis in $L^2(X,w(x)\,dx)$
constituted by eigenfunctions of $L$, that is, $L \varphi_n = \lambda_n \varphi_n$), is given by
$$
  R \varphi_n = \delta L^{-1/2} \varphi_n = \lambda_n^{-1/2} \delta \varphi_n
$$
and, for $f = \sum_{n=0}^{\infty} c_n \varphi_n \in L^2(X,w(x)\,dx)$,
$$
  Rf = \delta L^{-1/2} f = \sum_{n=0}^{\infty} c_n \lambda_n^{-1/2} \delta \varphi_n.
$$

Now, let us take
$$
  M = \delta \delta^*
$$
so $M\delta = \delta\delta^*\delta = \delta L$. Besides, % remembering that
% $\{\varphi_n\}_{n=0}^{\infty}$ is a orthonormal basis in $L^2(X,w(x)\,dx)$
% constituted by eigenfunctions of $L$, that is $L \varphi_n = \lambda_n \varphi_n$,
let us define
$$
  \psi_n = \delta \varphi_n,
$$
that satisfies $\delta^* \psi_n = \delta^* \delta \varphi_n = L \varphi_n = \lambda_n\varphi_n$.
Then
$$
  M \psi_n = \delta \delta^* \delta \varphi_n = \delta L \varphi_n = \delta \lambda_n\varphi_n
  = \lambda_n \delta \varphi_n = \lambda_n \psi_n,
$$
so $\{\psi_n\}_{n=0}^{\infty}$ are the eigenvalues of~$M$;
in particular, the operators $L$ and $M$ have the same spectrum.
Moreover, it is also easy to show (at least formally) that
$\{\psi_n\}_{n=0}^{\infty}$ is an orthonormal basis in $\delta(L^2(X,w(x)\,dx))$.

From the eigenfunctions $\{\psi_n\}_{n=0}^{\infty}$, we can, as with $L$, use the
spectral method to define many operators related to~$M$.
For instance, the semigroup $e^{-tM}$.
An important property is that, for ``any expression'' $\Phi(L)$, we have
\begin{equation}
\label{eq:Phi}
  \Phi(M) \delta = \delta \Phi(L).
\end{equation}

For any eigenvalue $\varphi_n$ it is clear that
$$
  \frac{\partial^2}{\partial t^2} \big( e^{-t\sqrt{M}}\delta L^{-1/2} \varphi_n)
  = \lambda_n^{1/2} e^{-t\sqrt{\lambda_n}} \psi_n
  = M e^{-t\sqrt{M}}\delta L^{-1/2} \varphi_n
$$
so, for $f \in L^2(X,w(x)\,dx)$, the function
\begin{equation}
\label{eq:Poissonconj}
  v(x,t) = e^{-t\sqrt{M}}\delta L^{-1/2}f(x)
\end{equation}
satisfies
\[
  \Bigg(\frac{\partial^2}{\partial t^2} - M \Bigg) v(x,t) = 0,
\]
that is similar to~\eqref{eq:harmonic}.
Actually, $v$ plays the role of the conjugate harmonic of~$u$, because
there is something similar to the Cauchy-Riemann equations;
indeed, by applying~\eqref{eq:Phi} and cancelling $\sqrt{L}$ with $L^{-1/2}$, we have
\begin{align*}
  \frac{\partial}{\partial t} v(x,t)
  &= \frac{\partial}{\partial t} e^{-t\sqrt{M}}\delta L^{-1/2}f(x)
  = -\sqrt{M} e^{-t\sqrt{M}}\delta L^{-1/2}f(x) \\
  &= -\delta \sqrt{L} e^{-t\sqrt{L}} L^{-1/2}f(x)
  = -\delta e^{-t\sqrt{L}} f(x)
  = -\delta u(x,t)
\end{align*}
and
\begin{align*}
  \delta^* v(x,t)
  &= \delta^* e^{-t\sqrt{M}}\delta L^{-1/2}f(x)
  = \delta^* \delta e^{-t\sqrt{L}} L^{-1/2}f(x)
  = L e^{-t\sqrt{L}} L^{-1/2}f(x) \\
  &= L^{1/2} e^{-t\sqrt{L}} f(x)
  = - \frac{\partial}{\partial t} e^{-t\sqrt{L}} f(x)
  = - \frac{\partial}{\partial t} u(x,t).
\end{align*}
By this reason, \eqref{eq:Poissonconj} is called the conjugate Poisson operator.
Finally, it is interesting to note that, in~\eqref{eq:Poisson}, instead of
$\lim_{t\to0} u(x,t) = f(x)$, we have, for~$v$,
$$
  \lim_{t\to0} v(x,t) = v(x,0) = \delta L^{-1/2} f(x) = Rf(x),
$$
i.e., we are getting the Riesz transform of~$f$.

All this theory that generalizes the classical harmonic analysis
requires, in the first place, to find the decomposition~\eqref{eq:decomposition}.
In practice, we also want to find decompositions such that, if we start with a ``good''
system $\{\varphi_n\}_{n=0}^{\infty}$ and we take $\psi_n = \delta\varphi_n$, then
the new system $\{\psi_n\}_{n=0}^{\infty}$ is ``good'' again.

Here, the easiest trick is to notice that the operator $L$
can be chosen ``except for an aditive constant'' because, if we add a constant,
the only effect is that all the eigenvalues are increased by this constant.
This can serve to find suitable decompositions
$$
  L - c = \delta^* \delta,
$$
where the constant $c$ must be chosen fulfilling $c<\lambda_0$ so that the
new eigenvalues $\{\lambda_n-c\}_{n=0}^\infty$ to be positive.
Sometimes, the useful decomposition to operate is done via two summands
$$
 L - c = \frac12(\delta^* \delta+\delta \delta^*),
$$
like for the harmonic oscillator (here $c=0$)
$$
  -\frac{\partial^2}{\partial x^2} + x^2
  = \frac12 \Big( {-\frac{\partial}{\partial x}} + x \Big)
    \Big( \frac{\partial}{\partial x} + x \Big)
  + \frac12 \Big( \frac{\partial}{\partial x} + x \Big)
    \Big( {-\frac{\partial}{\partial x}} + x \Big);
$$
this decomposition is the usual for expansions with respect to
the orthogonal system of Hermite functions.
% (See~\cite{StT}.)

Many different techniques and tricks have been followed for different operators,
and consequently there does not exist a standard way of defining the Riesz transform.
We can see many of them in~\cite{NSt}, a survey devoted to studying the $L^2$-theory
of Riesz transforms for many orthogonal expansions.

In the rest of the paper we are going to show an orthogonal system related to a second
order differential operator that is not easy to analyze with this scheme.
Actually, we have not been able to make a definition of the Riesz transform valid
in the corresponding $L^2$ space.

%%%%%%%%%%%%%%%%%%%%%%%%%%%%%%%%%%%%%%%%%%%%%%%%%%%%%%
\section{The case of Fourier-Neumann expansions}
%%%%%%%%%%%%%%%%%%%%%%%%%%%%%%%%%%%%%%%%%%%%%%%%%%%%%%

Let $J_\nu$ stand for the Bessel function of the first kind of order~$\nu$
(see \cite{Wat}, the well-known Watson's treatise on Bessel functions).
For $\alpha>-1$, the formula
\begin{equation*}
  \int_0^\infty J_{\alpha+2n+1}(x)J_{\alpha+2m+1}(x)\frac{dx}{x}
  = \frac{\delta_{nm}}{2(\alpha+2n+1)},
  \quad n,m = 0,1,2,\dots,
\end{equation*}
provides an orthonormal system $\{j_n^\alpha\}_{n=0}^\infty$ in
$L^2((0,\infty),x^{2\alpha+1}\,dx)$
% ($L^2(\mu_\alpha)$ from now on)
given by
\begin{equation*}
  j_n^\alpha(x) = \sqrt{2(\alpha+2n+1)}\,J_{\alpha+2n+1}(x) x^{-\alpha-1},
  \qquad n = 0,1,2,\dots\,.
\end{equation*}
For a function $f$, provided that the coefficients
\begin{equation*}
  c_k^\alpha(f) = \int_0^\infty f(y)j_k^\alpha(y) y^{2\alpha+1}\,dy,
  \qquad k = 0,1,2,\dots\,,
\end{equation*}
exist, consider its partial sums
\begin{equation*}
  S_n^\alpha(f,x) = \sum_{k=0}^n c_k^\alpha(f) j_k^\alpha(x),
  \qquad n = 0,1,2,\dots\,.
\end{equation*}
Series of the form $\sum_{n\ge0} a_n J_{\alpha+n}$ are usually
called Neumann series, hence we refer to $\sum_{k=0}^\infty c_k^\alpha(f)
j_k^\alpha(x)$ as to a Fourier-Neumann series.
The mean and almost everywhere convergence of these series have been
studied in~\cite{V,CiGPV,CiStV,CiV}.

It is well-known that the Bessel function $J_\nu(x)$ satisfies the differential equation
$$
  x^2 J''_\nu(x) + x J'_\nu(x) + (x^2-\nu^2) J_\nu(x) = 0.
$$
Then, it is easy to check that, if we denote by $L_\alpha$ the differential operator 
% (considered up to a constant, as seen above)
$$
  L_\alpha f \equiv \left(x^2\, \frac{d^2}{dx^2}
    + (2\alpha+3)x\, \frac{d}{dx} + x^2 + (\alpha+1)^2 \right)(f),
$$
we have
$$
  L_\alpha j_n^\alpha = (\alpha+2n+1)^2 j_n^\alpha,
$$
i.e., $j_n^\alpha$, $n = 0, 1, 2, \dots$, are eigenfunctions of $L_\alpha$ with eigenvalues $\lambda_n = (\alpha+2n+1)^2$. (Let us note we are seeing which differential equation is satisfied by the function $j_n^\alpha$, and it is a simple change of variable on the abovementioned differential equation for the Bessel functions; even more, the precise equation that is satisfied by~$j_n^\alpha$ can be found in~\cite[\S\,5.73, p.~158]{Wat}.)

The system $\{j_n^\alpha\}_{n=0}^{\infty}$ is not a basis on $L^2((0,\infty),x^{2\alpha+1}\,dx)$, but a basis in a proper subset $B_2 \varsubsetneq L^2((0,\infty),x^{2\alpha+1}\,dx)$ (with the same norm), that makes it be an interesting and peculiar system. Although it is not important in this moment, let us say by completeness that the space $B_2$ is formed by the functions $f \in L^2((0,\infty),x^{2\alpha+1}\,dx)$ such that the so-called ``modified'' Hankel transform of order~$\alpha$,
\begin{equation}
\label{eq:htr}
  H_\alpha(f)(x) = \int_0^{\infty} \frac{J_{\alpha}(xy)}{(xy)^{\alpha}} f(y) y^{2\alpha+1}\,dy,\quad x>0,
\end{equation}
is supported on $[0,1]$; see \cite{V,CiGPV,CiStV,CiV} for details, where not only is $B_2$ identified, but also
%$$
%  \overline{\spa\{j_n^{\alpha}\}_{n=0}^{\infty}} = B_p 
% \varsubsetneq L^p((0,\infty),x^{2\alpha+1}\,dx)
%$$
%(closure on $L^p((0,\infty),x^{2\alpha+1}\,dx)$) 
$$
  B_p = \overline{\spa\{j_n^{\alpha}\}_{n=0}^{\infty}} 
 \quad \text{(closure on $L^p((0,\infty),x^{2\alpha+1}\,dx)$)}
$$
for some range of $p$'s.

The density of $\spa\{j_n^{\alpha}\}_{n=0}^{\infty}$ in a proper subset of the whole $L^2$ (or~$L^p$) is not a handicap to use the spectral method of defining operators related to~$L_\alpha$; the only restriction is that we define these operators only for functions on the subset. In~\cite{BCiMPTV}, we have recently studied the heat and Poisson semigroups related to the differential operator~$L_\alpha$, as well as its fractional integrals. But, how to define the Riesz transform?

In this paper we are going to check that it does not seem to be an easy task, and that unexpected problems arise when trying to do the decomposition of $L_\alpha$ into two adjoint differential operators.

%%%%%%%%%%%%%%%%%%%%%%%%%%%%%%%%%%%%%%%%%%%%%%%%%%%%%%
\subsection{The decomposition of the operator}
%%%%%%%%%%%%%%%%%%%%%%%%%%%%%%%%%%%%%%%%%%%%%%%%%%%%%%

Let us decompose
$$
  L_\alpha \equiv x^2\, \frac{d^2}{dx^2}
    + (2\alpha+3)x\, \frac{d}{dx} + x^2 + (\alpha+1)^2 = \delta^* \delta,
$$
with $\delta$ and $-\delta^*$ two adjoint differential opperators, i.e., that verify
\begin{equation}
\label{eq:adjdd}
  \int_0^\infty f(x) \delta g(x) x^{2\alpha+1}\,dx
  = -\int_0^\infty g(x) \delta^*f(x) x^{2\alpha+1}\,dx
\end{equation}
for functions $f$ and $g$ good enough. 
%% YA LO DIGO M\'AS ABAJO:
% (in particular, that $x^{2\alpha+1}f(x)g(x)$ ``vanishes enough'' at $0$ and $\infty$).
%
The reason for taking this kind of decomposition instead of 
\eqref{eq:decomposition}--\eqref{eq:adjoint} is that the coefficient
of $\frac{d^2}{dx^2}$ in $L_\alpha$ (with positive eigenvalues) is $x^2$,
which is positive (otherwise, we need to use complex coefficients in 
$\delta$ and~$\delta^*$).
%
In what follows, all the derivatives involved in the differential operators $L_\alpha$, $\delta$ and~$\delta^*$, are with respect to the variable~$x$.

If we have two adjoint operators $\delta$ and $\delta^*$, it is easy to check whether $L_\alpha = \delta^* \delta$ or not. But the contrary is not true; that is, given $L_\alpha$, it is not easy to find adjoint operators $\delta$ and $\delta^*$ such that $L_\alpha = \delta^* \delta$. Moreover, there are many other variants that can be useful, such as $L_\alpha = \delta^* \delta + c$
%$L_\alpha = \frac12 \left(\delta^* \delta + \delta \delta^* \right)$,
or $L_\alpha - c = \frac12 \left(\delta^* \delta + \delta \delta^* \right)$,
for a constant~$c$, as was seen above.

The best way is to use a good computer algebra system such as Mathematica 
(but with its cons, see~\cite{CiV-M});
any other should be useful for the same purposes.
Although we will not reproduce the Mathematica code, all the formulas and decompositions that we show in what follows have been found with the help of this program.
Since $L_\alpha = x^2 \frac{d^2}{dx^2} + \cdots$, it seems reasonable to find decompositions with the form $\delta = x \frac{d}{dx} + \cdots$ and also $\delta^* = x \frac{d}{dx} + \cdots$. We have also tried operators with the more general form $\delta = x r(x) \frac{d}{dx} + \cdots$ and $\delta^* = \frac{x}{r(x)} \frac{d}{dx} + \cdots$ (for an arbitrary function $r(x)$) but nothing interesting seems to arise, so we do not show any of them in this paper. 

The first idea is to look for with Mathematica functions $p(x)$ and $q(x)$ such that
\begin{align*}
  \delta f = x \frac{df}{dx} + q(x) f, \\
  \delta^* f = x \frac{df}{dx} + p(x) f
\end{align*}
satisfy $L_\alpha = \delta^* \delta$ and \eqref{eq:adjdd} for functions $f,g$ that ``vanished enough'' in $0$ and~$\infty$ (making integrations by parts, this means that $f(x)g(x)x^{2\alpha+1}$ vanishes both in $x=0$ and when $x \to \infty$).

In this way we have found the following decomposition:
\begin{equation}
\label{eq:J}
\begin{aligned}
  \delta f = x \frac{df}{dx} + \left(\alpha+1 + \frac{xJ_1(x)}{J_0(x)} \right) f, \\
  \delta^* f = x \frac{df}{dx} + \left(\alpha+1 - \frac{xJ_1(x)}{J_0(x)} \right) f
\end{aligned}
\end{equation}
(we will call it the ``$J$-decomposition'').
This kind of decompositions $\delta^* \delta$ is very strange, because the infinity quantity of zeros of $J_0(x)$ makes the fraction $xJ_1(x)/J_0(x)$ be neither continuous nor derivable in infinitely many points, and it is not a function in~$L^p$ ($p \ge 1$).
% (Note: The inverse decomposition does not works: $\delta \delta^* \ne L_\alpha$.)

Moreover, if, as usual, $Y_\nu(x)$ denotes the Bessel function of the second kind and order~$\nu$, also the following decomposition appears:
\begin{equation}
\label{eq:Y}
\begin{aligned}
  \delta f = x \frac{df}{dx} + \left( \alpha+1 + \frac{xY_1(x)}{Y_0(x)} \right) f, \\
  \delta^* f = x \frac{df}{dx} + \left( \alpha+1 - \frac{xY_1(x)}{Y_0(x)} \right) f
\end{aligned}
\end{equation}
(we will call it the ``$Y$-decomposition''). It is also easy to get new decompositions by combining the $J$-decomposition~\eqref{eq:J} and the $Y$-decomposition~\eqref{eq:Y} by means of a process of mediation\footnote{The mediation is a very common process in diophantine approximation due to the property
$$
  \frac{a}{b} < \frac{c}{d} \text{\; (with denominators $b,d>0$)}
  \;\Rightarrow\; \frac{a}{b} < \frac{a+c}{b+d} < \frac{c}{d}; 
$$
it is used, for instance, in the definition of Farey and Brocot sequences.
It is funny to see that it also appears in the context of this paper.} of fractions (i.e., passing from $a/b$ and $c/d$ to $(a+c)/(b+d)$). In this way, we have the decomposition
\begin{equation}
\label{eq:med}
\begin{aligned}
  \delta f = x \frac{df}{dx} + \frac12 \left( \alpha+1 + \frac{ x J_1(x) 
    + x Y_1(x)}{J_0(x)+Y_0(x)} \right) f, \\
  \delta^* f = x \frac{df}{dx} + \frac12 \left( \alpha+1 - \frac{ x J_1(x) 
    + x Y_1(x)}{J_0(x)+Y_0(x)} \right) f.
\end{aligned}
\end{equation}
Even more, we can add two extra real parameters $r$ and $s$ and take
%\begin{equation}
%\label{eq:medrs}
%\begin{aligned}
%  \delta f = x \frac{df}{dx} + \left( \frac{r (\alpha J_0(x)+J_0(x)+x J_1(x))
%    + s (\alpha Y_0(x)+Y_0(x)+x Y_1(x))}{r J_0(x)+sY_0(x)} \right) f, \\
%  \delta^* f = x \frac{df}{dx} + \left( \frac{r (\alpha J_0(x)+J_0(x)-x J_1(x)) 
%    + s (\alpha Y_0(x)+Y_0(x)-x Y_1(x))}{r J_0(x)+s Y_0(x)} \right) f.
%\end{aligned}
%\end{equation}
\begin{equation}
\label{eq:medrs}
\begin{aligned}
  \delta f = x \frac{df}{dx} + \left( \alpha+1 + \frac{r x J_1(x)
    + s x Y_1(x)}{r J_0(x)+sY_0(x)} \right) f, \\
  \delta^* f = x \frac{df}{dx} + \left( \alpha+1 - \frac{r x J_1(x) 
    + s x Y_1(x)}{r J_0(x)+s Y_0(x)} \right) f.
\end{aligned}
\end{equation}
But we always get $J_\nu$ or $Y_\nu$ Bessel functions in the denominator, with infinitely many zeros. Recall that, for $x$ large enough, $J_\nu(x)$ and $Y_\nu(x)$ oscillate as a cosine or a sine respectively; more precisely, they satisfy the asymptotic formulas
\begin{align*}
  J_\nu(x) &= \left(\frac{2}{\pi x}\right)^{1/2}\left[ \cos\left(x-\frac{\nu\pi}2
    - \frac\pi4 \right) + O(x^{-1}) \right], \qquad x \to \infty, \\
  Y_\nu(x) &= \left(\frac{2}{\pi x}\right)^{1/2}\left[ \sin\left(x-\frac{\nu\pi}2
    - \frac\pi4 \right) + O(x^{-1}) \right], \qquad x \to \infty.
\end{align*}

By definition, the Riesz transform is $R = \delta (L_\alpha)^{-1/2}$, with $(L_\alpha)^{-1/2}$ being the fractional integral. Thus, if we apply this to $f = \sum_{n=0}^\infty c_n j_n^\alpha$, we get
\begin{align*}
  Rf(x) &= R\left(\sum_{n=0}^\infty c_n j_n^\alpha\right)(x)
  = \delta (L_\alpha)^{-1/2} \left(\sum_{n=0}^\infty c_n j_n^\alpha\right)(x) \\
  &= \delta \left(\sum_{n=0}^\infty
    \frac{c_n}{((\alpha+2n+1)^2)^{1/2}} \,j_n^\alpha\right)(x)
  = \sum_{n=0}^\infty \frac{c_n}{\alpha+2n+1} \, \delta j_n^\alpha (x).
\end{align*}
But (unlike other cases in which the Riesz transform is studied) it does not seem easy to find a ``good'' expression for $\delta j_n^\alpha (x)$, that in the case of~\eqref{eq:J} would be
\begin{multline*}
  \qquad
  \delta j_n^\alpha (x) = \sqrt{2(\alpha+2n+1)}\,x \frac{d}{dx} 
    \frac{J_{\alpha+2n+1}(x)}{x^{\alpha+1}} \\*
  + \sqrt{2(\alpha+2n+1)} \left(\alpha+1 + \frac{xJ_1(x)}{J_0(x)} \right) 
    \frac{J_{\alpha+2n+1}(x)}{x^{\alpha+1}},
  \qquad
\end{multline*}
so nothing ``smart'' seems to appear. (We will not insist on it, but the same can be said for the decompositions as in \eqref{eq:c}--\eqref{eq:complex} that we will see below, and also for~\eqref{eq:sqrt} and~\eqref{eq:complex-2}.)

%% ES LO MISMO QUE ANTES. NO APORTA NADA PONERLO.
%	Another thing that can be tried is to take $R f = \delta (L_\alpha)^{-1/2} f$ 
%	by means of the subordination formula for $(L_\alpha)^{-1/2}$, that is,
%	taking something like
%	$$
%	  Rf(x) = x \frac{d}{dx}\int_0^\infty (e^{-t\sqrt{L_\alpha}}f)(x)\, dt
%	  + \left(\alpha+1 + \frac{xJ_1(x)}{J_0(x)} \right)
%	    \int_0^\infty (e^{-t\sqrt{L_\alpha}}f)(x)\, dt.
%	$$
%	But, in the same way, we can not find in which way this expression may be useful.

A different idea is to use a constant $c$ and decompose
$$
  L_\alpha - c \equiv x^2\, \frac{d^2}{dx^2}
    + (2\alpha+3)x\, \frac{d}{dx} + x^2 + (\alpha+1)^2 - c = \delta^* \delta.
$$
We have found that this can be done with
\begin{equation}
\label{eq:c}
\begin{aligned}
  \delta f &= x \frac{df}{dx} + \left( 1+\alpha
    + \frac{x(J_{1+\sqrt{c}}(x)-J_{-1+\sqrt{c}}(x))}{2J_{\sqrt{c}}(x)} \right) f \\*
  &= x \frac{df}{dx} + \biggl( 1+\alpha
    - \frac{xJ'_{\sqrt{c}}(x)}{J_{\sqrt{c}}(x)} \biggr) f, \\
  \delta^* f &= x \frac{df}{dx} + \left( 1+\alpha
    - \frac{x(J_{1+\sqrt{c}}(x)-J_{-1+\sqrt{c}}(x))}{2J_{\sqrt{c}}(x)} \right) f \\*
  &= x \frac{df}{dx} + \biggl( 1+\alpha
    + \frac{xJ'_{\sqrt{c}}(x)}{J_{\sqrt{c}}(x)} \biggr) f.
\end{aligned}
\end{equation}
Of course, similar $Y$-decomposition can be found, as well as combinations between the $J$-decomposition and the $Y$-decomposition by means of the mediation process as in~\eqref{eq:med} and~\eqref{eq:medrs}. We do not see marvelous improvements, but it is interesting to note that, with $c=1/4$, we get trigonometric expressions;
and with $c=(\alpha+1)^2$, the Bessel functions in the decomposition depend on $\alpha$.
(For us, it sounds strange that, except for the summand $1+\alpha$, the decomposition can be taken independent of~$\alpha$.) 

An example of decomposition with $c=1/4$ is
\begin{equation}
\label{eq:tan}
\begin{aligned}
  \delta f &= x \frac{df}{dx} + \left( \frac32 +\alpha
    + x \tan(x) \right) f, \\
  \delta^* f &= x \frac{df}{dx} + \left( \frac12 +\alpha
    - x \tan(x) \right) f.
\end{aligned}
\end{equation}
Another (also with $c=1/4$):
\begin{equation}
\label{eq:cot}
\begin{aligned}
  \delta f &= x \frac{df}{dx} + \left( \frac32 +\alpha
    - x \cot(x) \right) f, \\
  \delta^* f &= x \frac{df}{dx} + \left( \frac12 +\alpha
    + x \cot(x) \right) f.
\end{aligned}
\end{equation}
In general (for $c=1/4$), and taking into account a mediation process as in~\eqref{eq:medrs}, we get
\begin{equation}
\label{eq:sincos}
\begin{aligned}
  \delta f &= x \frac{df}{dx} + \left( \frac{((3+2\alpha)s-2rx)\cos(x)
    + ((3+2\alpha)r+2sx)\sin(x)}{2s\cos(x)+2r\sin(x)} \right) f, \\
  \delta^* f &= x \frac{df}{dx} + \left( \frac{((1+2\alpha)s+2rx)\cos(x)
    + ((1+2\alpha)r-2sx)\sin(x)}{2s\cos(x)+2r\sin(x)} \right) f
\end{aligned}
\end{equation}
for arbitrary constants $r$ and~$s$. At least formally, these constants can be complex numbers;
in this way, taking $r=i$ and $s=1$, we get
\begin{equation}
\label{eq:complex}
\begin{aligned}
  \delta f &= x \frac{df}{dx} + \left( \frac32 +\alpha + ix \right) f, \\
  \delta^* f &= x \frac{df}{dx} + \left( \frac12 +\alpha - ix \right) f.
\end{aligned}
\end{equation}
But it seems difficult to apply these decompositions in the context of real functions.

%%%%%%%%%%%%%%%%%%%%%%%%%%%%%%%%%%%%%%%%%%%%%%%%%%%%%%
\subsection{Decompositions with two summands}
%%%%%%%%%%%%%%%%%%%%%%%%%%%%%%%%%%%%%%%%%%%%%%%%%%%%%%

Decompositions of the form
\begin{equation*}
%\label{eq:2sum}
    L_\alpha - c = \frac12 \left(\delta^* \delta + \delta \delta^* \right)
\end{equation*}
(with $\delta$ and $-\delta^*$ adjoint) can be also found. This time, completely
different expressions appear in the decomposition 
(without Bessel functions in a denominator!), 
but unfortunately nothing useful seems to arise.

For instance, we have checked that~ the above decomposition with two summands can be obtained with
\begin{equation}
\label{eq:sqrt}
\begin{aligned}
  \delta f &= x \frac{df}{dx} + \left( 1 +\alpha + \sqrt{c-x^2} \right) f, \\
  \delta^* f &= x \frac{df}{dx} + \left( 1 +\alpha - \sqrt{c-x^2} \right) f.
\end{aligned}
\end{equation}
But $x \in (0,\infty)$ so, now, we always get complex functions in a range of~$x$.
In particular, with $c=0$ we get
\begin{equation}
\label{eq:complex-2}
\begin{aligned}
  \delta f &= x \frac{df}{dx} + \left( 1 +\alpha + i x \right) f, \\
  \delta^* f &= x \frac{df}{dx} + \left( 1 +\alpha - i x \right) f.
\end{aligned}
\end{equation}
% Of course, it seems difficult to apply these decompositions in the context of real functions.

%%%%%%%%%%%%%%%%%%%%%%%%%%%%%%%%%%%%%%%%%%%%%%%%%%%%%%
\subsection{The ``non-modified'' case}
%%%%%%%%%%%%%%%%%%%%%%%%%%%%%%%%%%%%%%%%%%%%%%%%%%%%%%

%\begin{rem*}
We can wonder what would happen if we carry out the study of Riesz transforms in the Fourier-Neumann expansions setting in another way. The case studied up to now concerns the measure $x^{2\alpha+1}\,dx$, which is usually known as the ``modified'' case (because it is related with the modified Hankel transform~\eqref{eq:htr}), and that is used in~\cite{BCiMPTV, CiGPV, CiStV, CiV, V}. But we might be also interested in the case that deals with Lebesgue measure, or the ``non-modified'' case (that is related to the so-called non-modified Hankel transform, a simple variation of~\eqref{eq:htr}).

In this context, we consider the space $L^2((0,\infty),dx)$ in which the system $\{\phi_n^\alpha\}_{n=0}^{\infty}$, given by
\begin{equation*}
  \phi_n^\alpha(x) = \sqrt{2(\alpha+2n+1)}\,J_{\alpha+2n+1}(x) x^{-1/2},
  \qquad n = 0,1,2,\dots
\end{equation*}
is orthonormal. The associated second order differential operator is
$$
  \mathcal{L}_\alpha f \equiv \left(x^2\, \frac{d^2}{dx^2}
    + 2x\, \frac{d}{dx} + x^2 + \frac14 \right)(f),
$$
and we have
$$
  \mathcal{L}_\alpha \phi_n^\alpha = (\alpha+2n+1)^2 \phi_n^\alpha,
$$
i.e., $\phi_n^\alpha$, $n = 0, 1, 2, \dots$, are eigenfunctions of $\mathcal{L}_\alpha$ with eigenvalues $\lambda_n = (\alpha+2n+1)^2$. 

Again with the help of Mathematica, we have also found the corresponding expressions for the non-modified case, but as one can expect (we are making simple change of functions), nothing new appears. For instance, we get $\mathcal{L}_\alpha = \delta^* \delta$ with
\begin{align*}
  \delta f &= x \frac{df}{dx} + \left(\frac12 + \frac{xJ_1(x)}{J_0(x)} \right) f, \\
  \delta^* f &= x \frac{df}{dx} + \left(\frac12 - \frac{xJ_1(x)}{J_0(x)} \right) f,
\end{align*}
or, analogously, with its corresponding $Y$-decomposition.
Of course, the same can be said about the mediation process, the variation of the operator by an additive constant~$c$, as in $\mathcal{L}_\alpha - c = \delta^*\delta$ (one more time, trigonometric expressions appear for $c=1/4$), and the decompositions of the kind $\mathcal{L}_\alpha = \frac12 (\delta^* \delta + \delta \delta^*)$ or $\mathcal{L}_\alpha - c = \frac12 (\delta^* \delta + \delta \delta^*)$.
%\end{rem*}


%%%%%%%%%%%%%%%%%%%%%%%%%%%%%%%%%%%%%%%%%%%%%%%%%%%%%%
\begin{thebibliography}{10}
%%%%%%%%%%%%%%%%%%%%%%%%%%%%%%%%%%%%%%%%%%%%%%%%%%%%%%

\frenchspacing

%\bibitem{BMTU} \textsc{D. Buraczewski, T. Mart\'{\i}nez, J. L. Torrea, and R. Urban},
%{On the Riesz transform associated with ultraspherical polynomials},
%\textit{J. Anal. Math.} (to appear),
%\texttt{http://www.uam.es/joseluis.torrea/research.html}

%\bibitem{BaE} \textsc{D. Bakry and M. \'Emery},
%{Diffusions hypercontractives},
%S\'eminaire de probabilit�s (Strasbourg), tome 19,
%1983/84, p.~177--206, Lecture Notes in Math. \textbf{1123},
%Springer, Berlin, 1985.

\bibitem{BCiMPTV} \textsc{J. J. Betancor, \'O. Ciaurri, T. Mart\'{\i}nez, M. P\'erez, J. L. Torrea, and J. L. Varona},
{Heat and Poisson semigroups for Fourier-Neumann expansions},
\textit{Semigroup Forum}
\textbf{73} (2006), 129--142.

\bibitem{CiGPV} \textsc{\'O. Ciaurri, J. J. Guadalupe, M. P\'erez, and J. L. Varona},
{Mean and almost everywhere convergence of Fourier-Neumann series},
\textit{J. Math. Anal. Appl.}
\textbf{236} (1999), 125--147.

\bibitem{CiStV} \textsc{\'O. Ciaurri, K. Stempak, and J. L. Varona},
Mean Ces\`aro-type summability of Fourier-Neumann series,
\textit{Studia Sci. Math. Hung.}
\textbf{42} (2005), 413--430.

\bibitem{CiV-M} \textsc{\'O. Ciaurri and J. L. Varona},
{?`Podemos fiarnos de los c\'alculos efectuados con ordenador?},
%% How reliable are computer calculations? (Spanish),
\textit{Gac. R. Soc. Mat. Esp.}
\textbf{9} (2006), 483--514.

\bibitem{CiV} \textsc{\'O. Ciaurri and J. L. Varona},
{The surprising almost everywhere convergence of Fourier-Neumann series},
\textit{J. Comput. Appl. Math.} 
\textbf{233} (2009), 663--666. 

%\bibitem{GCMST} \textsc{J. Garc\'{\i}a-Cuerva, G. Mauceri, S. Meda, P. Sj\"{o}gren, and J. L. Torrea},
%{Functional calculus for the Ornstein-Uhlenbeck operator},
%\textit{J. Funct. Anal.} 
%\textbf{183} (2001), 413--450.

%\bibitem{HTVi} \textsc{E. Harboure, J. L. Torrea, and B. E. Viviani},
%{Riesz transforms for Laguerre expansions},
%\textit{Indiana Univ. Math. J.} 
%\textbf{55} (2006), 999--1014.

\bibitem{MST} \textsc{R. A. Mac\'{\i}as, C. Segovia, and J. L. Torrea},
{Heat diffusion maximal operators for Laguerre semigroups with negative parameters},
\textit{J. Funct. Anal.}
\textbf{229} (2005), 300--316.

\bibitem{Mu} \textsc{B. Muckenhoupt},
{Poisson integrals for Hermite and Laguerre expansions},
\textit{Trans. Amer. Math. Soc.}
\textbf{139} (1969), 231--242.

\bibitem{MuSt} \textsc{B. Muckenhoupt and E. M. Stein},
{Classical expansions and their relation to conjugate harmonic functions},
\textit{Trans. Amer. Math. Soc.}
\textbf{118} (1965), 17--92.

\bibitem{NSt} \textsc{A. Nowak and K. Stempak},
{$L^2$-theory of Riesz transforms for orthogonal expansions},
\textit{J. Fourier Anal. Appl.}
\textbf{12} (2006), 675--711.

\bibitem{Sj} \textsc{P. Sj\"ogren},
{Operators associated with the Hermite semigroup---a survey},
% Proceedings of the conference dedicated to Professor Miguel de Guzm\'an (El Escorial, 1996)
\textit{J. Fourier Anal. Appl.}
\textbf{3} (1997), Special Issue, 813--823.

% El libro rojo de Stein
\bibitem{Stein} \textsc{E. M. Stein},
\textit{Topics in harmonic analysis related to the Littlewood-Paley theory},
{Annals of Mathematics Studies} \textbf{63},
Princeton Univ. Press,
Princeton, NJ, 1970.

%\bibitem{Stein-2} \textsc{E. M. Stein},
%\textit{Singular integrals and differentiability properties of functions},
%Princeton Mathematical Series \textbf{30},
%Princeton Univ. Press,
%Princeton, NJ, 1970.

\bibitem{StT} \textsc{K. Stempak and J. L. Torrea},
{Poisson integrals and Riesz transforms for Hermite function expansions with weights},
\textit{J. Funct. Anal.}
\textbf{202} (2003), 443--472.

\bibitem{Th1} \textsc{S. Thangavelu},
{Riesz transforms and the wave equation for the Hermite operator},
\textit{Comm. Partial Differential Equations}
\textbf{15} (1990), 1199--1215.

\bibitem{Th2} \textsc{S. Thangavelu}, 
\textit{Lectures on Hermite and Laguerre expansions},
Princeton Univ. Press, Princeton, New Jersey, 1993.

\bibitem{Th3} \textsc{S. Thangavelu},
{On conjugate Poisson integrals and Riesz transforms for the Hermite expansions},
\textit{Colloq. Math.}
\textbf{64} (1993), 103--113.

\bibitem{T} \textsc{J. L. Torrea},
{Algunas observaciones sobre el semigrupo de Laguerre},
%% {Some remarks about the Laguerre semigroup} (Spanish), 
in
\textit{Margarita mathematica en memoria de Jos\'e Javier (Chicho) Guadalupe Hern\'andez},
p.~365--373, Universidad de La Rioja, Logro\~no, 2001.

\bibitem{V} \textsc{J. L. Varona},
{Fourier series of functions whose Hankel transform is supported on $[0,1]$},
\textit{Constr. Approx.}
\textbf{10} (1994), 65--75.

\bibitem{Wat} \textsc{G. N. Watson},
\textit{A Treatise on the Theory of Bessel Functions} (2nd edition),
Cambridge Univ. Press,
Cambridge, 1944.

\end{thebibliography}

\selectlanguage{spanish} % <--- Afecta a los cortes de palabra en las direcciones

\end{document}

